% quickderivs-doc.tex
%
% Copyright (C) 2026 James Petersen <m@jamespetersen.ca>
% Licensed under MIT. See LICENSE.

\documentclass{article}
\usepackage[utf8]{inputenc}

\usepackage[upright]{quickderivs}
\dvsetnormal
\usepackage[margin=1.5in]{geometry}
\usepackage{verbatim}

\title{Derivatives using the {\tt quickderivs} package}
\author{James Petersen}
\date{March 25th, 2026 (Version 0.1.1)}

\begin{document}

\maketitle

This package contains commands to typeset mathematical derivatives more
efficiently.

\section{Options}
\begin{itemize}
	\item \verb|upright|: enables support for upright derivative symbols,
		such as
		\makeatletter\(\quickderivs@unslant{\partial}\)\makeatother,
		as compared to \(\partial\). These can be toggled
		using the \verb|\dvsetupright| and \verb|\dvsetnormal| commands.
		If \verb|upright| is passed, then the style is set
		to upright by default.

	\item \verb|partial|: the \verb|\dv| and \verb|\df| commands now
		default partial derivative symbols. The last star argument
		now switches to normal derivative symbols.
\end{itemize}

\section{Commands}
\begin{itemize}
	\item	\verb|\dv[*][|\(\langle\)\textit{upper value}\(\rangle\)%
		\verb|][|\(\langle\)\textit{upper power}\(\rangle\)\verb|]{|%
		\(\langle\)\textit{lower variable(s)}\(\rangle\)\verb|}[|%
		\(\langle\)\textit{differential symbol}\(\rangle\)\verb|][*]|

		This is the main command of the package. It typesets derivatives in general.
		The only mandatory parameter is the lower variables. With a single
		variable, it typesets a normal derivative: \verb|\dv{x}| becomes
		\[
			\dv x.
		\]

		The differential symbol parameter changes the differential symbol,
		with the following supported parameters:
		\verb|d D e E p|. They produce the following outputs.
		\begin{center}
			\begin{tabular}{cccccc}
				&\verb|d|&\verb|D|&\verb|e|&\verb|E|&\verb|p|\\[1em]
				\dvsetupright
				upright
				&\(\displaystyle\dv x[d]\)
				&\(\displaystyle\dv x[D]\)
				&\(\displaystyle\dv x[e]\)
				&\(\displaystyle\dv x[E]\)
				&\(\displaystyle\dv x[p]\)\\[2em]
				normal
				\dvsetnormal
				&\(\displaystyle\dv x[d]\)
				&\(\displaystyle\dv x[D]\)
				&\(\displaystyle\dv x[e]\)
				&\(\displaystyle\dv x[E]\)
				&\(\displaystyle\dv x[p]\)
			\end{tabular}
		\end{center}
		Additional symbol can be added by defining the command
		\verb|\quickderivs@chrs@|\(\langle\)\textit{new char}\(\rangle\)
		to the desired symbol and
		\verb|\quickderivs@chrsspac@|\(\langle\)\textit{new char}\(\rangle\)
		to the spacing this character should have between it and a superscript.

		The last star parameter toggles between regular and partial derivatives.
		It takes priority over the differential symbol parameter.
		Setting the \verb|partial| parameter of the package makes
		derivatives default to partial derivatives, and switch to regular
		derivatives with the star.

		The lower variable(s) parameter can take many variables in a row,
		delimited by commas: \verb|\dv{x,y}*| becomes
		\[
			\dv{x,y}*.
		\]
		Simply specifying a number will make that the power of the preceding
		variable: \verb|\dv{x2}| becomes
		\[
			\dv{x2}.
		\]
		Numbers can be used as separators between multiple variables as well:
		\verb|\dv{x2y3}*| becomes
		\[
			\dv{x2y3}*.
		\]
		To instead have the number as a part of the derivative, or any
		other complex command, enclose it in square brackets:
		\verb|\dv{{(2x)},{\bar{y}}}*| becomes
		\[
			\dv{{(2x)},{\bar y}}*.
		\]
		To specify a nonstandard superscript, use a semicolon:
		\verb|\dv{x;r,y;s}*| becomes
		\[
			\dv{x;r,y;s}*.
		\]

		The upper value is self-explanatory: \verb|\dv[f]{x}| becomes
		\[
			\dv[f]x.
		\]
		The upper power is also simple. Expanding on an earlier example,
		\verb|\dv[][r+s]{x;r,y;s}*| becomes
		\[
			\dv[][r+s]{x;r,y;s}*.
		\]
		Finally, the first star writes the fractions inline:
		\verb|\dv*[f]{x}| becomes \(\dv*[f]x\).

	\item \verb|\df[*]{|\(\langle\)\textit{lower variable(s)}\(\rangle\)%
		\verb|}[|\(\langle\)\textit{differential symbol}\verb|][*]|

			The latter three parameters are identical to those in the \verb|\dv|
			command. The first star controls the spacing between the differentials:

			\begin{center}
				\begin{tabular}{cc}
					with star&without star\\
					\df{x,y,z}&\df*{x,y,z}
				\end{tabular}
			\end{center}

		\item \verb|\dvsetupright| and \verb|\dvsetnormal|

			These commands switch between upright and normal characters for the package.
			\verb|\dvsetupright| is only defined if the \verb|upright| parameter
			is supplied to the package.
\end{itemize}

\end{document}
